
\documentclass{article}
\usepackage{colortbl}
\usepackage{makecell}
\usepackage{multirow}
\usepackage{supertabular}

\begin{document}

\newcounter{utterance}

\centering \large Interaction Transcript for game `dond', experiment `semi\_en', episode 10 with Aya{-}Expanse{-}32B{-}t0.0.
\vspace{24pt}

{ \footnotesize  \setcounter{utterance}{1}
\setlength{\tabcolsep}{0pt}
\begin{supertabular}{c@{$\;$}|p{.15\linewidth}@{}p{.15\linewidth}p{.15\linewidth}p{.15\linewidth}p{.15\linewidth}p{.15\linewidth}}
    \# & $\;$A & \multicolumn{4}{c}{Game Master} & $\;\:$B\\
    \hline

    \theutterance \stepcounter{utterance}  
    & & \multicolumn{4}{p{0.6\linewidth}}{
        \cellcolor[rgb]{0.9,0.9,0.9}{
            \makecell[{{p{\linewidth}}}]{
                \texttt{\tiny{[P1$\langle$GM]}}
                \texttt{You are playing a negotiation game in which you have to agree on how to divide a set of items among you and another player.} \\
\\ 
\texttt{Rules:} \\
\texttt{(a) You and the other player are given a set of items. Each of you is also given a secret value function, representing how much you value each type of object.} \\
\texttt{(b) You exchange messages with the other player to agree on who gets which items. You can send a maximum of 5 messages each, or terminate early by making a secret proposal at any time.} \\
\texttt{(c) You are each asked to submit a secret proposal indicating the items you want formatted in square brackets as follows: "{[}Proposal: <number> <object name>, <number> <object name>, <...>{]}"} \\
\texttt{(d) If your proposals are complementary, i.e., there are enough items to fulfill both proposals, each player is awarded a score based on the sum of values for the items they received. Otherwise, both of you get zero points.} \\
\texttt{(e) Your goal is to maximize the score you receive.} \\
\\ 
\texttt{Let us start.} \\
\\ 
\texttt{The set of available items is:} \\
\\ 
\texttt{2 remotes, 2 magnets, 4 watches} \\
\\ 
\texttt{Your secret value function is:} \\
\\ 
\texttt{remote: 3, magnet: 2, watch: 0} \\
\\ 
\texttt{IMPORTANT: Your messages, unless it is the secret proposal, are directly transmitted to the other player, so do not include any response to the rules or text announcing your message. To make a secret proposal, use the indicated format. Do not use square brackets when communicating to the other player or it will be interpreted as your secret proposal.} \\
            }
        }
    }
    & \\ \\

    \theutterance \stepcounter{utterance}  
    & \multicolumn{4}{p{0.6\linewidth}}{
        \cellcolor[rgb]{0.9,0.9,0.9}{
            \makecell[{{p{\linewidth}}}]{
                \texttt{\tiny{[P1$\rangle$GM]}}
                \texttt{Understanding the dynamics and keeping my secret value function in mind, I'd initiate the negotiation with the following message:} \\
\\ 
\texttt{"Hi, I think it's crucial we find a fair distribution that satisfies both of us. Given the items, I propose we each take 1 remote, as they seem to be in high demand. Then, we can divide the magnets and watches based on our preferences."} \\
\\ 
\texttt{This message sets a tone for cooperation and fairness, while also suggesting a baseline division that could work for both of us. Given my value function, I'm prioritizing the remotes, which are worth the most to me. By proposing a 1:1 split of these, I'm hoping the other player will consider it a fair starting point.} \\
\\ 
\texttt{After this message, I'd be interested in hearing their response and adjusting my strategy based on their reaction and subsequent messages. If they seem amenable to the remote split, I'd focus on trying to secure at least one magnet, as those are my next most valuable item. If they push back against the remote split, I might need to reevaluate and consider other strategies.} \\
\\ 
\texttt{Remember, I can send up to 5 messages or make a secret proposal at any time. For now, I'm aiming for an open and fair negotiation to maximize the chance of a mutually beneficial agreement.} \\
            }
        }
    }
    & & \\ \\

    \theutterance \stepcounter{utterance}  
    & & & \multicolumn{2}{p{0.3\linewidth}}{
        \cellcolor[rgb]{0.9,0.9,0.9}{
            \makecell[{{p{\linewidth}}}]{
                \texttt{\tiny{[GM$|$GM]}}
                \texttt{continue} \\
            }
        }
    }
    & & \\ \\

    \theutterance \stepcounter{utterance}  
    & & & \multicolumn{2}{p{0.3\linewidth}}{
        \cellcolor[rgb]{0.9,0.9,0.9}{
            \makecell[{{p{\linewidth}}}]{
                \texttt{\tiny{[GM$|$GM]}}
                \texttt{Understanding the dynamics and keeping my secret value function in mind, I'd initiate the negotiation with the following message:} \\
\\ 
\texttt{"Hi, I think it's crucial we find a fair distribution that satisfies both of us. Given the items, I propose we each take 1 remote, as they seem to be in high demand. Then, we can divide the magnets and watches based on our preferences."} \\
\\ 
\texttt{This message sets a tone for cooperation and fairness, while also suggesting a baseline division that could work for both of us. Given my value function, I'm prioritizing the remotes, which are worth the most to me. By proposing a 1:1 split of these, I'm hoping the other player will consider it a fair starting point.} \\
\\ 
\texttt{After this message, I'd be interested in hearing their response and adjusting my strategy based on their reaction and subsequent messages. If they seem amenable to the remote split, I'd focus on trying to secure at least one magnet, as those are my next most valuable item. If they push back against the remote split, I might need to reevaluate and consider other strategies.} \\
\\ 
\texttt{Remember, I can send up to 5 messages or make a secret proposal at any time. For now, I'm aiming for an open and fair negotiation to maximize the chance of a mutually beneficial agreement.} \\
            }
        }
    }
    & & \\ \\

    \theutterance \stepcounter{utterance}  
    & & \multicolumn{4}{p{0.6\linewidth}}{
        \cellcolor[rgb]{0.9,0.9,0.9}{
            \makecell[{{p{\linewidth}}}]{
                \texttt{\tiny{[GM$\rangle$P2]}}
                \texttt{You are playing a negotiation game in which you have to agree on how to divide a set of items among you and another player.} \\
\\ 
\texttt{Rules:} \\
\texttt{(a) You and the other player are given a set of items. Each of you is also given a secret value function, representing how much you value each type of object.} \\
\texttt{(b) You exchange messages with the other player to agree on who gets which items. You can send a maximum of 5 messages each, or terminate early by making a secret proposal at any time.} \\
\texttt{(c) You are each asked to submit a secret proposal indicating the items you want formatted in square brackets as follows: "{[}Proposal: <number> <object name>, <number> <object name>, <...>{]}"} \\
\texttt{(d) If your proposals are complementary, i.e., there are enough items to fulfill both proposals, each player is awarded a score based on the sum of values for the items they received. Otherwise, both of you get zero points.} \\
\texttt{(e) Your goal is to maximize the score you receive.} \\
\\ 
\texttt{Let us start.} \\
\\ 
\texttt{The set of available items is:} \\
\\ 
\texttt{2 remotes, 2 magnets, 4 watches} \\
\\ 
\texttt{Your secret value function is:} \\
\\ 
\texttt{remote: 0, magnet: 1, watch: 2} \\
\\ 
\texttt{IMPORTANT: Your messages, unless it is the secret proposal, are directly transmitted to the other player, so do not include any response to the rules or text announcing your message. To make a secret proposal, use the indicated format. Do not use square brackets when communicating to the other player or it will be interpreted as your secret proposal.} \\
\\ 
\\ 
\texttt{Understanding the dynamics and keeping my secret value function in mind, I'd initiate the negotiation with the following message:} \\
\\ 
\texttt{"Hi, I think it's crucial we find a fair distribution that satisfies both of us. Given the items, I propose we each take 1 remote, as they seem to be in high demand. Then, we can divide the magnets and watches based on our preferences."} \\
\\ 
\texttt{This message sets a tone for cooperation and fairness, while also suggesting a baseline division that could work for both of us. Given my value function, I'm prioritizing the remotes, which are worth the most to me. By proposing a 1:1 split of these, I'm hoping the other player will consider it a fair starting point.} \\
\\ 
\texttt{After this message, I'd be interested in hearing their response and adjusting my strategy based on their reaction and subsequent messages. If they seem amenable to the remote split, I'd focus on trying to secure at least one magnet, as those are my next most valuable item. If they push back against the remote split, I might need to reevaluate and consider other strategies.} \\
\\ 
\texttt{Remember, I can send up to 5 messages or make a secret proposal at any time. For now, I'm aiming for an open and fair negotiation to maximize the chance of a mutually beneficial agreement.} \\
            }
        }
    }
    & \\ \\

    \theutterance \stepcounter{utterance}  
    & & & \multicolumn{4}{p{0.6\linewidth}}{
        \cellcolor[rgb]{0.9,0.9,0.9}{
            \makecell[{{p{\linewidth}}}]{
                \texttt{\tiny{[GM$\langle$P2]}}
                \texttt{**Message 1:**} \\
\\ 
\texttt{"Let's aim for a balanced agreement. I suggest we start by splitting the remotes evenly, so we each get one. This seems like a fair starting point, ensuring neither of us is left without this seemingly desirable item."} \\
\\ 
\texttt{**Strategy:**} \\
\\ 
\texttt{My initial proposal aims to establish a foundation for a fair deal. By suggesting a 1:1 split of the remotes, I'm addressing the potentially high demand for this item while also revealing my willingness to cooperate. Given my value function, I'm indeed prioritizing remotes, but this message doesn't give away my entire strategy.} \\
\\ 
\texttt{**Potential Follow{-}up:**} \\
\\ 
\texttt{{-} If the other player agrees to the remote split, I'll focus on magnets and watches. I might express my preference for magnets, given their value to me, and suggest a trade or division of these items.} \\
\texttt{{-} In case they disagree or counter{-}propose, I'll need to assess their value for remotes and adjust my strategy. I could offer a trade, e.g., both remotes for a higher number of watches, depending on their response.} \\
\\ 
\texttt{**Key Points:**} \\
\\ 
\texttt{{-} Keeping an open and friendly tone encourages collaboration.} \\
\texttt{{-} Suggesting a fair initial split sets a positive tone for the negotiation.} \\
\texttt{{-} My value function guides my priorities, but I'll adapt to the other player's responses to reach a mutually beneficial agreement.} \\
            }
        }
    }
     \\ \\

    \theutterance \stepcounter{utterance}  
    & & & \multicolumn{2}{p{0.3\linewidth}}{
        \cellcolor[rgb]{0.9,0.9,0.9}{
            \makecell[{{p{\linewidth}}}]{
                \texttt{\tiny{[GM$|$GM]}}
                \texttt{continue} \\
            }
        }
    }
    & & \\ \\

    \theutterance \stepcounter{utterance}  
    & & & \multicolumn{2}{p{0.3\linewidth}}{
        \cellcolor[rgb]{0.9,0.9,0.9}{
            \makecell[{{p{\linewidth}}}]{
                \texttt{\tiny{[GM$|$GM]}}
                \texttt{**Message 1:**} \\
\\ 
\texttt{"Let's aim for a balanced agreement. I suggest we start by splitting the remotes evenly, so we each get one. This seems like a fair starting point, ensuring neither of us is left without this seemingly desirable item."} \\
\\ 
\texttt{**Strategy:**} \\
\\ 
\texttt{My initial proposal aims to establish a foundation for a fair deal. By suggesting a 1:1 split of the remotes, I'm addressing the potentially high demand for this item while also revealing my willingness to cooperate. Given my value function, I'm indeed prioritizing remotes, but this message doesn't give away my entire strategy.} \\
\\ 
\texttt{**Potential Follow{-}up:**} \\
\\ 
\texttt{{-} If the other player agrees to the remote split, I'll focus on magnets and watches. I might express my preference for magnets, given their value to me, and suggest a trade or division of these items.} \\
\texttt{{-} In case they disagree or counter{-}propose, I'll need to assess their value for remotes and adjust my strategy. I could offer a trade, e.g., both remotes for a higher number of watches, depending on their response.} \\
\\ 
\texttt{**Key Points:**} \\
\\ 
\texttt{{-} Keeping an open and friendly tone encourages collaboration.} \\
\texttt{{-} Suggesting a fair initial split sets a positive tone for the negotiation.} \\
\texttt{{-} My value function guides my priorities, but I'll adapt to the other player's responses to reach a mutually beneficial agreement.} \\
            }
        }
    }
    & & \\ \\

    \theutterance \stepcounter{utterance}  
    & & \multicolumn{4}{p{0.6\linewidth}}{
        \cellcolor[rgb]{0.9,0.9,0.9}{
            \makecell[{{p{\linewidth}}}]{
                \texttt{\tiny{[P1$\langle$GM]}}
                \texttt{**Message 1:**} \\
\\ 
\texttt{"Let's aim for a balanced agreement. I suggest we start by splitting the remotes evenly, so we each get one. This seems like a fair starting point, ensuring neither of us is left without this seemingly desirable item."} \\
\\ 
\texttt{**Strategy:**} \\
\\ 
\texttt{My initial proposal aims to establish a foundation for a fair deal. By suggesting a 1:1 split of the remotes, I'm addressing the potentially high demand for this item while also revealing my willingness to cooperate. Given my value function, I'm indeed prioritizing remotes, but this message doesn't give away my entire strategy.} \\
\\ 
\texttt{**Potential Follow{-}up:**} \\
\\ 
\texttt{{-} If the other player agrees to the remote split, I'll focus on magnets and watches. I might express my preference for magnets, given their value to me, and suggest a trade or division of these items.} \\
\texttt{{-} In case they disagree or counter{-}propose, I'll need to assess their value for remotes and adjust my strategy. I could offer a trade, e.g., both remotes for a higher number of watches, depending on their response.} \\
\\ 
\texttt{**Key Points:**} \\
\\ 
\texttt{{-} Keeping an open and friendly tone encourages collaboration.} \\
\texttt{{-} Suggesting a fair initial split sets a positive tone for the negotiation.} \\
\texttt{{-} My value function guides my priorities, but I'll adapt to the other player's responses to reach a mutually beneficial agreement.} \\
            }
        }
    }
    & \\ \\

    \theutterance \stepcounter{utterance}  
    & \multicolumn{4}{p{0.6\linewidth}}{
        \cellcolor[rgb]{0.9,0.9,0.9}{
            \makecell[{{p{\linewidth}}}]{
                \texttt{\tiny{[P1$\rangle$GM]}}
                \texttt{**Message 2:**} \\
\\ 
\texttt{"Given our initial agreement on remotes, I'd like to discuss the magnets and watches. I have to admit, magnets hold a higher value for me. How about we aim for one magnet each and then divide the watches based on our preferences? I'm open to suggestions on how to fairly split the watches."} \\
\\ 
\texttt{**Strategy Update:**} \\
\\ 
\texttt{With the remote division settled (assuming their agreement), I'm now focusing on the next most valuable items for me—magnets. By proposing an even split of magnets, I'm again prioritizing fairness while ensuring I get something of higher value. This approach also allows us to divide the watches in a way that might work for both of us, depending on their preferences.} \\
\\ 
\texttt{**Contingency Planning:**} \\
\\ 
\texttt{{-} If they value magnets highly too, I might need to reconsider and suggest a trade{-}off. For instance, I could offer them both magnets in exchange for a larger share of watches or another item they might prefer.} \\
\texttt{{-} Should they agree to the magnet split, we can negotiate the watches. I'd aim for at least one watch, given their zero value to me, to ensure I get something in that category.} \\
\\ 
\texttt{**Negotiation Dynamics:**} \\
\\ 
\texttt{{-} Being transparent about my preferences for magnets while also offering a fair split sets a positive tone.} \\
\texttt{{-} Showing flexibility on the watches indicates a willingness to compromise, which could encourage a similar attitude from the other player.} \\
\texttt{{-} I'll continue to adapt my strategy based on their responses to reach a mutually beneficial agreement within the 5{-}message limit.} \\
            }
        }
    }
    & & \\ \\

    \theutterance \stepcounter{utterance}  
    & & & \multicolumn{2}{p{0.3\linewidth}}{
        \cellcolor[rgb]{0.9,0.9,0.9}{
            \makecell[{{p{\linewidth}}}]{
                \texttt{\tiny{[GM$|$GM]}}
                \texttt{continue} \\
            }
        }
    }
    & & \\ \\

    \theutterance \stepcounter{utterance}  
    & & & \multicolumn{2}{p{0.3\linewidth}}{
        \cellcolor[rgb]{0.9,0.9,0.9}{
            \makecell[{{p{\linewidth}}}]{
                \texttt{\tiny{[GM$|$GM]}}
                \texttt{**Message 2:**} \\
\\ 
\texttt{"Given our initial agreement on remotes, I'd like to discuss the magnets and watches. I have to admit, magnets hold a higher value for me. How about we aim for one magnet each and then divide the watches based on our preferences? I'm open to suggestions on how to fairly split the watches."} \\
\\ 
\texttt{**Strategy Update:**} \\
\\ 
\texttt{With the remote division settled (assuming their agreement), I'm now focusing on the next most valuable items for me—magnets. By proposing an even split of magnets, I'm again prioritizing fairness while ensuring I get something of higher value. This approach also allows us to divide the watches in a way that might work for both of us, depending on their preferences.} \\
\\ 
\texttt{**Contingency Planning:**} \\
\\ 
\texttt{{-} If they value magnets highly too, I might need to reconsider and suggest a trade{-}off. For instance, I could offer them both magnets in exchange for a larger share of watches or another item they might prefer.} \\
\texttt{{-} Should they agree to the magnet split, we can negotiate the watches. I'd aim for at least one watch, given their zero value to me, to ensure I get something in that category.} \\
\\ 
\texttt{**Negotiation Dynamics:**} \\
\\ 
\texttt{{-} Being transparent about my preferences for magnets while also offering a fair split sets a positive tone.} \\
\texttt{{-} Showing flexibility on the watches indicates a willingness to compromise, which could encourage a similar attitude from the other player.} \\
\texttt{{-} I'll continue to adapt my strategy based on their responses to reach a mutually beneficial agreement within the 5{-}message limit.} \\
            }
        }
    }
    & & \\ \\

    \theutterance \stepcounter{utterance}  
    & & \multicolumn{4}{p{0.6\linewidth}}{
        \cellcolor[rgb]{0.9,0.9,0.9}{
            \makecell[{{p{\linewidth}}}]{
                \texttt{\tiny{[GM$\rangle$P2]}}
                \texttt{**Message 2:**} \\
\\ 
\texttt{"Given our initial agreement on remotes, I'd like to discuss the magnets and watches. I have to admit, magnets hold a higher value for me. How about we aim for one magnet each and then divide the watches based on our preferences? I'm open to suggestions on how to fairly split the watches."} \\
\\ 
\texttt{**Strategy Update:**} \\
\\ 
\texttt{With the remote division settled (assuming their agreement), I'm now focusing on the next most valuable items for me—magnets. By proposing an even split of magnets, I'm again prioritizing fairness while ensuring I get something of higher value. This approach also allows us to divide the watches in a way that might work for both of us, depending on their preferences.} \\
\\ 
\texttt{**Contingency Planning:**} \\
\\ 
\texttt{{-} If they value magnets highly too, I might need to reconsider and suggest a trade{-}off. For instance, I could offer them both magnets in exchange for a larger share of watches or another item they might prefer.} \\
\texttt{{-} Should they agree to the magnet split, we can negotiate the watches. I'd aim for at least one watch, given their zero value to me, to ensure I get something in that category.} \\
\\ 
\texttt{**Negotiation Dynamics:**} \\
\\ 
\texttt{{-} Being transparent about my preferences for magnets while also offering a fair split sets a positive tone.} \\
\texttt{{-} Showing flexibility on the watches indicates a willingness to compromise, which could encourage a similar attitude from the other player.} \\
\texttt{{-} I'll continue to adapt my strategy based on their responses to reach a mutually beneficial agreement within the 5{-}message limit.} \\
            }
        }
    }
    & \\ \\

    \theutterance \stepcounter{utterance}  
    & & & \multicolumn{4}{p{0.6\linewidth}}{
        \cellcolor[rgb]{0.9,0.9,0.9}{
            \makecell[{{p{\linewidth}}}]{
                \texttt{\tiny{[GM$\langle$P2]}}
                \texttt{**Message 3:**} \\
\\ 
\texttt{"I understand the importance of magnets, and I'm glad we could agree on a fair distribution of remotes. Now, regarding the watches, I'd like to propose a simple solution. Since they hold less value for me, I'm content with receiving just one watch. You can have the remaining three. This way, we ensure a quick and fair resolution to this part of the negotiation."} \\
\\ 
\texttt{**Tactical Decision:**} \\
\\ 
\texttt{By offering to take only one watch, I'm making a strategic concession. Given my value function, this move might seem counterintuitive, but it serves multiple purposes:} \\
\\ 
\texttt{{-} **Speed and Simplicity:** This proposal simplifies the negotiation, allowing us to reach an agreement faster. With one less item category to discuss, we can save time and potentially avoid complex trade{-}offs.} \\
\\ 
\texttt{{-} **Demonstrating Flexibility:** Showing willingness to accept less than my maximum value demonstrates flexibility and a genuine desire to reach a mutually agreeable outcome. This could encourage the other player to reciprocate.} \\
\\ 
\texttt{{-} **Focus on Higher{-}Value Items:** By securing one magnet and one watch, I'm ensuring I get something in each category I value. This leaves me with a satisfactory outcome, even with a smaller watch count.} \\
\\ 
\texttt{**Potential Outcomes:**} \\
\\ 
\texttt{{-} If the other player agrees, we have a clear and fair division: {[}1 remote, 1 magnet, 1 watch{]} for me and {[}1 remote, 1 magnet, 3 watches{]} for them.} \\
\texttt{{-} In case they counter{-}propose, I'll evaluate their offer. I might consider trading one of my magnets for an additional watch, depending on their proposal.} \\
\\ 
\texttt{**Negotiation Progress:**} \\
\\ 
\texttt{{-} Three messages down, two to go (if needed). I'm aiming for a clean and efficient negotiation, maximizing my score while respecting the other player's interests.} \\
\texttt{{-} My strategy remains adaptable, ready to adjust based on their responses to reach a successful conclusion.} \\
            }
        }
    }
     \\ \\

    \theutterance \stepcounter{utterance}  
    & & & \multicolumn{2}{p{0.3\linewidth}}{
        \cellcolor[rgb]{0.9,0.9,0.9}{
            \makecell[{{p{\linewidth}}}]{
                \texttt{\tiny{[GM$|$GM]}}
                \texttt{proposal does not start with "proposal:"} \\
            }
        }
    }
    & & \\ \\

\end{supertabular}
}

\end{document}
